\documentclass[11pt]{article}

\usepackage[utf8]{inputenc}
\usepackage[T1]{fontenc}
\usepackage[spanish,es-nodecimaldot]{babel}
\usepackage{amsmath,amssymb,mathtools}
\usepackage{geometry}
\usepackage{booktabs}
\usepackage{array}
\usepackage{enumitem}
\usepackage{hyperref}
\usepackage{xcolor}
\usepackage{siunitx}
\usepackage{float}
\usepackage{graphicx}

\geometry{margin=2.5cm}
\hypersetup{
    colorlinks=true,
    linkcolor=blue!50!black,
    urlcolor=blue!50!black
}

\newcommand{\figplaceholder}[3]{
\begin{figure}[H]
\centering
\fbox{
\begin{minipage}[c][6.0cm][c]{0.82\textwidth}
\centering
\textbf{Espacio reservado para la figura}\\[4pt]
#3
\end{minipage}}
\caption{#1}
\label{#2}
\end{figure}
}

\title{Plataforma Antivibratoria Activa para un Sensor de Ultrasonido\\Documento Final Ajustado al Enunciado}
\author{Base final para Overleaf}
\date{21 de abril de 2026}

\begin{document}

\maketitle
\tableofcontents
\newpage

\section{Introduccion}

En el laboratorio se realizan pruebas sobre papas utilizando un sensor de ultrasonido. Durante la medicion, las vibraciones de la mesa o base donde esta montado el sensor alteran su posicion y producen errores en la distancia medida. En consecuencia, el problema se formula como un problema de \textbf{rechazo de perturbaciones}: se desea que la plataforma donde se encuentra el sensor permanezca lo mas inmovil posible aunque la base vibre.

El sistema propuesto es SISO y su planta nominal es de al menos tercer orden. La variable manipulada es el voltaje aplicado al actuador, mientras que la salida de interes es la posicion de la plataforma del sensor. El trabajo se organiza siguiendo exactamente las partes pedidas en el enunciado: modelado realista, obtencion del modelo matematico, comparacion entre modelos, dise\~no del controlador y evaluacion del desempe\~no.

\section{Contexto de aplicacion}

El sistema estudiado en este proyecto se plantea pensando en una implementacion de laboratorio orientada a la deteccion de plagas en papas. En ese montaje experimental, la deteccion se realiza mediante un sensor de ultrasonido. Este tipo de sensor es especialmente sensible a perturbaciones mecanicas, porque cualquier vibracion de la estructura donde esta montado puede alterar la distancia medida y producir errores en la lectura.

El sensor se ubica sobre un montaje instalado en una mesa de laboratorio. Si la mesa sufre vibraciones o pequenas perturbaciones externas, la posicion del sensor cambia respecto del punto de medicion, aun cuando el objeto inspeccionado no se haya movido realmente. Por eso, el problema no debe entenderse como uno de seguimiento de trayectoria, sino principalmente como uno de \textbf{regulacion alrededor de un punto fijo}.

En este trabajo, el objetivo de control se define imponiendo que el \textit{setpoint} del sistema sea siempre cero:
\[
r(t)=0.
\]
Esto significa que se desea mantener constante la posicion del sensor respecto de su equilibrio nominal. En otras palabras, el controlador debe actuar de manera que la salida del sistema permanezca lo mas cercana posible a cero cuando aparezcan perturbaciones sobre la mesa o sobre la base del montaje.

Desde el punto de vista fisico, mantener el \textit{setpoint} en cero significa evitar que la plataforma donde esta montado el sensor cambie de posicion ante vibraciones del entorno. Esa es precisamente la razon de modelar el sistema como una plataforma antivibratoria activa: el resorte y el amortiguador aportan aislamiento pasivo, mientras que el actuador aporta una fuerza de correccion para reducir aun mas el movimiento transmitido al sensor.

\section{1. Modelado realista del sistema}

\subsection{Descripcion fisica del sistema}

El sistema realista se compone de una base inferior que vibra por efecto del entorno del laboratorio, una plataforma superior sobre la cual se monta el sensor de ultrasonido, un resorte que une la plataforma con la base, un amortiguador viscoso que disipa energia y un actuador electromecanico que aplica una fuerza de correccion para contrarrestar el efecto de las vibraciones.

La idea fisica es clara: si la base vibra, la plataforma tendera a moverse. El resorte y el amortiguador aportan aislamiento pasivo, mientras que el actuador permite introducir una accion activa de control. Por eso el sistema se interpreta como una plataforma antivibratoria activa.

\subsection{Variables del modelo}

Se definen las siguientes variables:

\begin{center}
\small
\begin{tabular}{>{\raggedright\arraybackslash}p{2.7cm} >{\raggedright\arraybackslash}p{9.5cm}}
\toprule
\textbf{Simbolo} & \textbf{Interpretacion} \\
\midrule
$x(t)$ & posicion absoluta de la plataforma donde esta montado el sensor \\
$z(t)$ & posicion absoluta de la base \\
$q(t)=x(t)-z(t)$ & desplazamiento relativo entre plataforma y base \\
$u(t)$ & voltaje aplicado al actuador \\
$i(t)$ & corriente interna del actuador \\
$F_a(t)$ & fuerza generada por el actuador \\
$y(t)=x(t)$ & salida regulada del sistema \\
\bottomrule
\end{tabular}
\end{center}

\subsection{Aspectos realistas incluidos}

El modelo realista implementado en Simulink/Simscape debe representar la fisica del sistema y, ademas, incluir efectos que no esten presentes en el modelo matematico aproximado. En este proyecto, el efecto adicional mas importante del modelo realista es la \textbf{saturacion del actuador}, ya que en la practica el actuador no puede generar un voltaje o una fuerza ilimitada. Si se desea enriquecer aun mas la simulacion, tambien pueden incluirse restricciones mecanicas adicionales o limites de recorrido.

\figplaceholder{Modelo realista implementado en Simulink/Simscape.}{fig:modelo_realista}{Reemplazar este espacio por la captura del modelo realista donde se identifiquen claramente la base, la plataforma, el resorte, el amortiguador, el actuador y el punto de montaje del sensor.}

\section{2. Obtencion del modelo matematico}

\subsection{Modelo fisico completo}

Se introduce la coordenada relativa
\begin{equation}
q(t)=x(t)-z(t),
\end{equation}
porque el resorte y el amortiguador dependen de la deformacion y de la velocidad relativa entre sus extremos.

Aplicando la segunda ley de Newton a la plataforma:
\begin{equation}
m\ddot x = -k(x-z)-c(\dot x-\dot z)+F_a.
\end{equation}

Usando la relacion $x=q+z$, se obtiene:
\begin{equation}
m\ddot q + c\dot q + kq = F_a - m\ddot z.
\label{eq:mecanica_completa}
\end{equation}

Para el actuador se usa un modelo electromecanico lineal:
\begin{equation}
F_a = K_f i,
\label{eq:fuerza_corriente}
\end{equation}
\begin{equation}
L\dot i + Ri + K_e\dot q = u.
\label{eq:ecuacion_electrica}
\end{equation}

Sustituyendo \eqref{eq:fuerza_corriente} en \eqref{eq:mecanica_completa}, el modelo fisico completo queda:
\begin{equation}
m\ddot q + c\dot q + kq = K_f i - m\ddot z,
\label{eq:modelo_completo_1}
\end{equation}
\begin{equation}
L\dot i + Ri + K_e\dot q = u.
\label{eq:modelo_completo_2}
\end{equation}

Estas ecuaciones representan el sistema fisico completo. La perturbacion de base aparece de forma explicita en el termino $-m\ddot z$.

\subsection{Planta nominal para analisis y dise\~no}

Para construir la planta nominal de control, la perturbacion de base se anula temporalmente:
\[
z(t)=0,\qquad \dot z(t)=0,\qquad \ddot z(t)=0.
\]

En ese caso, \eqref{eq:modelo_completo_1}--\eqref{eq:modelo_completo_2} quedan:
\begin{equation}
m\ddot q + c\dot q + kq = K_f i,
\end{equation}
\begin{equation}
L\dot i + Ri + K_e\dot q = u.
\end{equation}

Los estados se definen como:
\[
x_1=q,\qquad x_2=\dot q,\qquad x_3=i.
\]

Con esta eleccion, la planta nominal se escribe como:
\begin{equation}
\dot x = Ax+Bu,
\end{equation}
\begin{equation}
y = Cx+Du,
\end{equation}
con
\begin{equation}
A=
\begin{bmatrix}
0 & 1 & 0\\
-\dfrac{k}{m} & -\dfrac{c}{m} & \dfrac{K_f}{m}\\
0 & -\dfrac{K_e}{L} & -\dfrac{R}{L}
\end{bmatrix},
\qquad
B=
\begin{bmatrix}
0\\
0\\
\dfrac{1}{L}
\end{bmatrix},
\end{equation}
\begin{equation}
C=
\begin{bmatrix}
1 & 0 & 0
\end{bmatrix},
\qquad
D=0.
\end{equation}

La observacion importante es que la planta nominal tiene \textbf{solo una entrada}, que es $u(t)$. La perturbacion de base no se incorpora como una segunda entrada del modelo nominal cuando se construyen $A$, $B$, $C$ y $D$.

\subsection{Funcion de transferencia nominal}

La funcion de transferencia nominal se obtiene con:
\begin{equation}
G(s)=C(sI-A)^{-1}B+D.
\end{equation}

Desarrollando esta expresion:
\begin{equation}
G(s)=\frac{Q(s)}{U(s)}
=
\frac{K_f}{(Ls+R)(ms^2+cs+k)+K_fK_e s}.
\label{eq:tf_nominal}
\end{equation}

Expandiendo el denominador:
\begin{equation}
G(s)=
\frac{K_f}
{Lm\,s^3 + (Lc+Rm)s^2 + (Lk+Rc+K_fK_e)s + Rk}.
\end{equation}

Esto confirma que la planta nominal es de \textbf{tercer orden}. La razon es que se necesitan tres estados independientes para describir la dinamica: $q$, $\dot q$ e $i$.

\subsection{Analisis de estabilidad}

La ecuacion caracteristica de la planta nominal es:
\[
Lm\,s^3 + (Lc+Rm)s^2 + (Lk+Rc+K_fK_e)s + Rk = 0.
\]

Si se define
\[
a_3=Lm,\qquad a_2=Lc+Rm,\qquad a_1=Lk+Rc+K_fK_e,\qquad a_0=Rk,
\]
entonces el criterio de Routh-Hurwitz para un polinomio cubico exige:
\begin{equation}
a_3>0,\quad a_2>0,\quad a_1>0,\quad a_0>0,\quad a_2a_1>a_3a_0.
\end{equation}

La tabla de Routh correspondiente es:
\[
\begin{array}{c|cc}
s^3 & a_3 & a_1\\
s^2 & a_2 & a_0\\
s^1 & \dfrac{a_2a_1-a_3a_0}{a_2} & 0\\
s^0 & a_0 & 0
\end{array}
\]

La primera columna de esa tabla debe ser estrictamente positiva para que todos los polos queden en el semiplano izquierdo. Por eso aparece la desigualdad adicional
\[
a_2a_1>a_3a_0,
\]
que es la condicion no trivial del problema.

Como los parametros fisicos son positivos, la condicion principal es:
\begin{equation}
(Lc+Rm)(Lk+Rc+K_fK_e)>LmRk.
\end{equation}

Por tanto, la estabilidad se estudia sobre la planta nominal de tercer orden obtenida a partir del modelo fisico.

\subsubsection{Interpretacion fisica}

Cada coeficiente del polinomio caracteristico tiene una lectura fisica clara:

\begin{itemize}[leftmargin=1.2cm]
\item $a_3=Lm$ mezcla la inercia mecanica con la dinamica electrica del actuador;
\item $a_2=Lc+Rm$ combina amortiguamiento viscoso con disipacion electrica;
\item $a_1=Lk+Rc+K_fK_e$ incorpora rigidez, disipacion y acoplamiento electromecanico;
\item $a_0=Rk$ representa el efecto restaurador estatico ponderado por la resistencia electrica.
\end{itemize}

El producto $K_fK_e$ merece una observacion especial. Ese termino nace del hecho de que la corriente genera fuerza sobre la plataforma y, a la vez, el movimiento de la plataforma induce una fuerza contraelectromotriz en la parte electrica. Por eso el actuador no solo agrega un estado, sino que modifica el amortiguamiento efectivo del sistema.

\subsubsection{Implicaciones para el sistema}

La estabilidad de la planta nominal significa que, sin importar las condiciones iniciales, la respuesta libre decae con el tiempo. Sin embargo, eso no garantiza por si solo un buen rechazo de perturbaciones. Un sistema puede ser estable y aun asi vibrar demasiado, asentarse muy lento o transmitir movimiento suficiente para deteriorar la medicion del sensor.

La desigualdad
\[
(Lc+Rm)(Lk+Rc+K_fK_e)>LmRk
\]
permite leer varias tendencias:

\begin{itemize}[leftmargin=1.2cm]
\item si $c$ aumenta, el sistema tiende a ser mas amortiguado;
\item si $R$ aumenta, la corriente cambia menos agresivamente y se reduce el efecto de la dinamica electrica rapida;
\item si $K_fK_e$ aumenta, el acoplamiento electromecanico tiende a aportar amortiguamiento adicional;
\item si $L$ aumenta demasiado, el actuador responde con mas rezago y puede deteriorar la separacion entre el polo electrico y los polos mecanicos;
\item si la desigualdad se viola, aparece al menos un polo en el semiplano derecho y la planta nominal se vuelve inestable.
\end{itemize}

En consecuencia, la estabilidad abierta de la planta de tercer orden es una condicion necesaria para continuar con el dise\~no, pero no suficiente para afirmar que la plataforma sera adecuada para el sensor. La evaluacion final debe mirar tambien amortiguamiento, rapidez y vibracion residual.

\section{3. Comparacion entre modelos}

El enunciado pide comparar el modelo realista implementado en Simulink/Simscape con el modelo matematico aproximado en lazo abierto frente a una entrada escalon. El objetivo de esta comparacion es verificar que el modelo teorico capture la dinamica dominante del sistema y, al mismo tiempo, mostrar las diferencias que aparecen cuando se consideran efectos fisicos adicionales en el modelo realista.

\subsection{Lazo abierto teorico con funcion de transferencia}

\figplaceholder{Respuesta en lazo abierto del modelo teorico obtenido a partir de la funcion de transferencia.}{fig:lazo_abierto_tf}{Insertar aqui la grafica de la respuesta al escalon del modelo teorico.}

En esta figura se deben discutir el tiempo de subida, el tiempo de establecimiento, el sobreimpulso y el valor final del modelo aproximado.

\subsection{Lazo abierto Simscape}

\figplaceholder{Respuesta en lazo abierto del modelo realista en Simscape.}{fig:lazo_abierto_simscape}{Insertar aqui la grafica de la respuesta al escalon del modelo realista.}

La comparacion entre las dos respuestas permite identificar diferencias en rapidez, amortiguamiento, saturacion o cualquier otro efecto que el modelo realista capture y que el modelo teorico no incluya.

\section{4. Dise\~no del controlador}

\subsection{Modelo reducido para sintonia IMC-PID}

Para la sintonia del controlador se utiliza un modelo reducido de segundo orden. Este paso no cambia el hecho de que la planta nominal completa sea de tercer orden; simplemente se adopta una aproximacion mas simple para obtener un PID ideal en forma cerrada.

\subsubsection{Por que se reduce el modelo}

El modelo completo es de tercer orden porque incluye dos estados mecanicos y un estado electrico. Esa es la descripcion apropiada para la planta nominal real. Sin embargo, un PID por IMC se obtiene de manera mucho mas clara cuando la planta de dise\~no es de segundo orden, ya que asi el controlador final puede escribirse directamente en forma ideal y su interpretacion fisica resulta inmediata.

La reduccion no significa que el tercer estado desaparezca fisicamente. Lo que se asume es que la dinamica electrica del actuador es suficientemente rapida frente a la dinamica mecanica dominante, de modo que puede aproximarse de forma cuasiestatica para fines de sintonia.

\subsubsection{Reduccion fisica desde la ecuacion electrica}

Partimos del modelo nominal completo sin perturbacion:
\[
m\ddot q + c\dot q + kq = K_f i,
\]
\[
L\dot i + Ri + K_e\dot q = u.
\]

Si la constante de tiempo electrica
\[
\tau_e=\frac{L}{R}
\]
es pequena frente a la escala de tiempo mecanica dominante, entonces en la banda de interes puede aproximarse
\[
L\dot i \approx 0.
\]

La ecuacion electrica queda:
\[
Ri + K_e\dot q \approx u,
\]
de donde
\begin{equation}
i \approx \frac{u-K_e\dot q}{R}.
\label{eq:i_aproximada}
\end{equation}

Sustituyendo \eqref{eq:i_aproximada} en la ecuacion mecanica:
\[
m\ddot q + c\dot q + kq = K_f\left(\frac{u-K_e\dot q}{R}\right).
\]

Agrupando terminos:
\begin{equation}
m\ddot q + \left(c+\frac{K_fK_e}{R}\right)\dot q + kq = \frac{K_f}{R}u.
\label{eq:modelo_reducido_fisico}
\end{equation}

La ecuacion \eqref{eq:modelo_reducido_fisico} ya es un modelo reducido de segundo orden con significado fisico claro. A partir de ella, la funcion de transferencia reducida mas fiel al sistema seria:
\begin{equation}
G_{r,f}(s)=\frac{Q(s)}{U(s)}\approx \frac{K_f/R}{ms^2+\left(c+\frac{K_fK_e}{R}\right)s+k}.
\label{eq:tf_reducida_fisica}
\end{equation}

Esta expresion muestra que al reducir:

\begin{itemize}[leftmargin=1.2cm]
\item se elimina el polo electrico rapido;
\item se conserva el comportamiento dominante de masa, amortiguamiento y rigidez;
\item el acoplamiento electromecanico no desaparece por completo, sino que queda absorbido en el amortiguamiento equivalente y en la ganancia de actuacion.
\end{itemize}

\subsubsection{Respuesta forzada del modelo reducido}

Hasta este punto se ha escrito el modelo reducido destacando la relacion entre entrada de control y salida. Sin embargo, para este proyecto no basta con estudiar la respuesta libre o la transferencia desde la accion de control. Como el objetivo principal es el rechazo de perturbaciones, tambien debe quedar explicita la \textbf{respuesta forzada}.

Para ello se introduce una fuerza de perturbacion equivalente $f_p(t)$ sobre la plataforma. Esta fuerza puede interpretarse de dos maneras:

\begin{itemize}[leftmargin=1.2cm]
\item como una fuerza externa efectiva transmitida por la mesa o por el montaje;
\item o como una representacion equivalente del efecto inercial de la base, de forma que
\[
f_p(t)\equiv -m\ddot z(t).
\]
\end{itemize}

Con esa notacion, el modelo reducido fisicamente consistente se escribe como:
\begin{equation}
m\ddot q + \left(c+\frac{K_fK_e}{R}\right)\dot q + kq = \frac{K_f}{R}u + f_p(t).
\label{eq:modelo_reducido_forzado}
\end{equation}

En Laplace:
\begin{equation}
Q(s)=\frac{K_f/R}{ms^2+\left(c+\frac{K_fK_e}{R}\right)s+k}U(s)
+\frac{1}{ms^2+\left(c+\frac{K_fK_e}{R}\right)s+k}F_p(s).
\label{eq:reducido_forzado_laplace}
\end{equation}

La ecuacion \eqref{eq:reducido_forzado_laplace} es importante porque separa claramente dos contribuciones:

\begin{itemize}[leftmargin=1.2cm]
\item la parte debida a la accion de control;
\item la parte debida a la perturbacion externa.
\end{itemize}

Por tanto, la salida total puede entenderse como la suma de una respuesta asociada al control y una respuesta forzada por la perturbacion. En el caso de regulacion con \textit{setpoint} fijo, esta segunda parte es precisamente la que se quiere minimizar.

Con los valores dados para la parte mecanica:
\[
m=1,\qquad c=5,\qquad k=100,
\]
se toma:
\begin{equation}
G_r(s)=\frac{1}{ms^2+cs+k}
=
\frac{1}{s^2+5s+100}.
\label{eq:planta_reducida}
\end{equation}

La expresion \eqref{eq:planta_reducida} es aun mas simple que \eqref{eq:tf_reducida_fisica}. Aqui se hizo una segunda simplificacion orientada a dise\~no: se tomaron directamente los parametros mecanicos principales y se absorbieron en la sintonia los detalles de ganancia del actuador y del amortiguamiento adicional. Esto no se hace porque la parte electrica sea irrelevante, sino porque el objetivo inmediato es obtener un PID interpretable y facil de implementar.

Si se quiere mantener explicita la perturbacion tambien en esta version simplificada, el modelo de dise\~no puede escribirse como:
\begin{equation}
m\ddot q + c\dot q + kq = u + f_p(t),
\label{eq:modelo_reducido_sintonia_forzado}
\end{equation}
de donde
\begin{equation}
Q(s)=\frac{1}{ms^2+cs+k}U(s)+\frac{1}{ms^2+cs+k}F_p(s).
\label{eq:gr_forzado}
\end{equation}

Es decir, en la forma mas simple del modelo reducido, tanto la accion de control como la fuerza de perturbacion excitan la misma dinamica mecanica de segundo orden. Esto es exactamente lo que hace visible la respuesta forzada en el modelo de sintonia.

\subsubsection{Que se conserva y que se pierde con esta reduccion}

La reduccion a \eqref{eq:planta_reducida} conserva la dinamica mecanica dominante, la frecuencia natural principal y la estructura basica de oscilador masa-resorte-amortiguador. Eso basta para construir una ley de control inicial.

Lo que se pierde es la fidelidad exacta del tercer orden:

\begin{itemize}[leftmargin=1.2cm]
\item el polo electrico ya no aparece explicitamente;
\item la fase del sistema a frecuencias mas altas puede diferir de la planta completa;
\item la agresividad del controlador calculado sobre el modelo reducido puede no coincidir exactamente con la respuesta del modelo realista.
\end{itemize}

Por esa razon, la reduccion se usa solo como herramienta de sintonia. La verificacion final del desempeno debe hacerse sobre la planta de tercer orden y sobre el modelo realista en Simscape.

\subsubsection{Lectura dinamica del modelo reducido}

Para
\[
G_r(s)=\frac{1}{s^2+5s+100},
\]
la frecuencia natural y el factor de amortiguamiento son:
\begin{equation}
\omega_n=\sqrt{\frac{k}{m}}=10\ \text{rad/s},
\qquad
\zeta=\frac{c}{2\sqrt{mk}}=\frac{5}{20}=0.25.
\end{equation}

Eso significa que el modelo reducido es estable pero subamortiguado. En consecuencia, el sistema tiende naturalmente a oscilar antes de asentarse, y por eso el controlador no solo debe estabilizar sino tambien aumentar el amortiguamiento efectivo de la plataforma.

Como modelo deseado de lazo cerrado se escoge:
\begin{equation}
F(s)=\frac{1}{\lambda s+1}.
\end{equation}

Esta notacion sigue la convencion utilizada en clase para IMC, donde el parametro de sintonia se denota por $\lambda$. En las notas de IMC, para rechazo a perturbaciones se propone un filtro de la forma
\[
F(s)=\frac{\lambda s+1}{(\lambda s+1)^n}.
\]
En este proyecto, al trabajar con una planta reducida de segundo orden y buscar un equivalente PID ideal con accion derivativa, la forma efectiva utilizada para la sintonia queda:
\[
F(s)=\frac{1}{\lambda s+1}.
\]
Por eso, el papel que antes se habia escrito con otro simbolo se deja aqui expresado directamente con la notacion del curso.

Si se identifica el filtro con la dinamica deseada de lazo cerrado, se obtiene:
\begin{equation}
T_d(s)=\frac{1}{\lambda s+1}.
\end{equation}

Aplicando la sintonia IMC:
\begin{equation}
C(s)=\frac{T_d(s)}{G_r(s)\left[1-T_d(s)\right]}
=
\frac{s^2+5s+100}{\lambda s}.
\label{eq:control_imc}
\end{equation}

En forma general, si
\[
G_r(s)=\frac{1}{ms^2+cs+k},
\]
entonces
\begin{equation}
C(s)=\frac{ms^2+cs+k}{\lambda s}
=
\frac{m}{\lambda}s+\frac{c}{\lambda}+\frac{k}{\lambda}\frac{1}{s}.
\label{eq:control_imc_general}
\end{equation}

\subsubsection{Lectura del lazo cerrado frente a perturbaciones}

Si se usa el modelo reducido forzado \eqref{eq:gr_forzado} y se cierra el lazo con realimentacion unitaria, entonces para el caso principal del proyecto, donde
\[
r(t)=0,
\]
la salida frente a una fuerza de perturbacion queda dada por:
\begin{equation}
\frac{Q(s)}{F_p(s)}=\frac{G_r(s)}{1+C(s)G_r(s)}.
\label{eq:disturbance_cl}
\end{equation}

La expresion \eqref{eq:disturbance_cl} muestra por que la respuesta forzada debe tenerse en cuenta: el controlador no solo cambia la respuesta al escalon de referencia, sino tambien la forma en que las perturbaciones se transmiten hacia la plataforma. En esta aplicacion, esta segunda funcion es la mas importante, porque el objetivo no es mover la mesa a distintos puntos, sino impedir que el sensor se desplace cuando aparece una excitacion externa.

\subsection{Forma PID ideal}

La forma ideal del PID es:
\begin{equation}
C_{PID}(s)=K_c\left(1+\frac{1}{\tau_I s}+\tau_D s\right).
\label{eq:pid_ideal}
\end{equation}

Expandiendo:
\begin{equation}
C_{PID}(s)=K_c+\frac{K_c}{\tau_I}\frac{1}{s}+K_c\tau_D s.
\label{eq:pid_expandido}
\end{equation}

Por otra parte, de \eqref{eq:control_imc} se obtiene:
\begin{equation}
C(s)=\frac{1}{\lambda}s+\frac{5}{\lambda}+\frac{100}{\lambda s}.
\label{eq:imc_expandido}
\end{equation}

Si primero se compara la forma general \eqref{eq:control_imc_general} con la expresion expandida del PID ideal \eqref{eq:pid_expandido}, se obtiene:
\begin{equation}
K_c=\frac{c}{\lambda},
\qquad
\tau_I=\frac{c}{k},
\qquad
\tau_D=\frac{m}{c}.
\label{eq:parametros_pid_generales}
\end{equation}

Es decir, antes de sustituir valores numericos ya puede verse que la ganancia proporcional depende del parametro de sintonia $\lambda$, mientras que los tiempos integral y derivativo quedan determinados por la relacion entre amortiguamiento, rigidez e inercia del modelo reducido.

Comparando \eqref{eq:pid_expandido} y \eqref{eq:imc_expandido} termino a termino:
\begin{equation}
K_c=\frac{5}{\lambda},
\end{equation}
\begin{equation}
\frac{K_c}{\tau_I}=\frac{100}{\lambda},
\end{equation}
\begin{equation}
K_c\tau_D=\frac{1}{\lambda}.
\end{equation}

De aqui se obtiene:
\begin{equation}
\tau_I=\frac{5}{100}=0.05\ \text{s},
\qquad
\tau_D=\frac{1}{5}=0.2\ \text{s}.
\end{equation}

\subsection{Mantener $P=50$}

Si se quiere conservar la ganancia proporcional en
\[
P=K_c=50,
\]
entonces:
\begin{equation}
50=\frac{5}{\lambda}
\quad\Rightarrow\quad
\lambda=0.1\ \text{s}.
\end{equation}

Esta eleccion es consistente con la logica de la sintonia IMC vista en clase: $\lambda$ es el parametro que ajusta la rapidez del lazo cerrado y, en general, se termina afinando por simulacion. En este trabajo, ademas de esa interpretacion, se impuso la condicion practica de mantener $P=50$, por lo cual el valor de $\lambda$ queda determinado de forma directa.

Por tanto, el controlador PID ideal final es:
\begin{equation}
\boxed{
C(s)=50\left(1+\frac{1}{0.05\,s}+0.2\,s\right)
}
\label{eq:pid_final}
\end{equation}

En forma paralela:
\[
K_p=50,\qquad K_i=1000,\qquad K_d=10.
\]

\subsection{Equivalencia con el bloque PID de Simulink}

Si se usa el bloque
\begin{equation}
C_b(s)=P\left(1+\frac{1}{Is}+D\frac{N}{1+N/s}\right),
\label{eq:bloque_pid}
\end{equation}
la equivalencia con el PID ideal debe hacerse con cuidado, porque el bloque \textbf{no} usa directamente los parametros $K_p$, $K_i$ y $K_d$ de la forma paralela. El bloque usa otra parametrizacion.

\subsubsection{Paso 1: partir del PID ideal ya obtenido}

El controlador ideal calculado en \eqref{eq:pid_final} es:
\begin{equation}
C_{PID}(s)=50\left(1+\frac{1}{0.05\,s}+0.2\,s\right).
\label{eq:pid_ideal_bloque}
\end{equation}

Si se expande \eqref{eq:pid_ideal_bloque}, se obtiene:
\begin{equation}
C_{PID}(s)=50+\frac{50}{0.05}\frac{1}{s}+50(0.2)s.
\end{equation}

Como
\[
\frac{50}{0.05}=1000,\qquad 50(0.2)=10,
\]
queda:
\begin{equation}
C_{PID}(s)=50+\frac{1000}{s}+10s.
\label{eq:pid_paralela}
\end{equation}

De aqui se lee la forma paralela:
\[
K_p=50,\qquad K_i=1000,\qquad K_d=10.
\]

\subsubsection{Paso 2: entender que el bloque no usa $K_i$ ni $K_d$}

La forma del bloque \eqref{eq:bloque_pid} es:
\[
C_b(s)=P\left(1+\frac{1}{Is}+D\frac{N}{1+N/s}\right).
\]

Si momentaneamente se ignora el filtro derivativo, la estructura basica del bloque es:
\[
P\left(1+\frac{1}{Is}+Ds\right).
\]

Comparando esta expresion con la forma ideal general:
\[
K_c\left(1+\frac{1}{\tau_I s}+\tau_D s\right),
\]
se observa que la correspondencia correcta es:
\begin{equation}
P=K_c,
\qquad
I=\tau_I,
\qquad
D=\tau_D.
\label{eq:equiv_ideal_bloque}
\end{equation}

Esto significa que:

\begin{itemize}[leftmargin=1.2cm]
\item el parametro \texttt{P} del bloque es la ganancia proporcional total;
\item el parametro \texttt{I} del bloque es el \textbf{tiempo integral}, no la ganancia integral;
\item el parametro \texttt{D} del bloque es el \textbf{tiempo derivativo}, no la ganancia derivativa.
\end{itemize}

Esa diferencia es importante porque muchas veces se confunden:
\[
I \neq K_i,
\qquad
D \neq K_d.
\]

En este caso:
\[
K_i=1000,\qquad K_d=10,
\]
pero los valores que debe recibir el bloque son:
\[
I=0.05,\qquad D=0.2.
\]

\subsubsection{Paso 3: sustitucion numerica directa}

Usando \eqref{eq:equiv_ideal_bloque} y los parametros del PID ideal obtenido:
\[
K_c=50,\qquad \tau_I=0.05,\qquad \tau_D=0.2,
\]
se obtiene directamente:
\begin{equation}
P=50,
\qquad
I=0.05,
\qquad
D=0.2.
\label{eq:PID_directos}
\end{equation}

\subsubsection{Paso 4: explicar el papel del filtro derivativo}

El termino derivativo del bloque no es $Ds$ puro, sino:
\begin{equation}
D\frac{N}{1+N/s}.
\label{eq:term_deriv_block}
\end{equation}

Para entenderlo mejor, se multiplica numerador y denominador por $s$:
\begin{equation}
D\frac{N}{1+N/s}
=
D\frac{Ns}{s+N}.
\label{eq:term_deriv_block_2}
\end{equation}

Esta expresion muestra que la derivada del bloque esta \textbf{filtrada}. Eso se hace para evitar que una derivada ideal pura amplifique demasiado el ruido de medicion.

Ahora bien, si $N$ es suficientemente grande en comparacion con la banda de interes, entonces:
\[
s+N \approx N,
\]
y por tanto:
\begin{equation}
\frac{Ns}{s+N}\approx \frac{Ns}{N}=s.
\end{equation}

En consecuencia,
\begin{equation}
D\frac{Ns}{s+N}\approx Ds.
\end{equation}

Esa es la razon por la cual el bloque puede aproximar el comportamiento del PID ideal, pero de una manera mas robusta frente a ruido.

\subsubsection{Paso 5: eleccion de $N$}

El parametro $N$ no sale de la comparacion con el PID ideal. No aparece en \eqref{eq:pid_final} porque es un parametro adicional del filtro derivativo del bloque.

Por eso, $N$ se elige como parametro de implementacion. Un valor inicial razonable es:
\[
N=10.
\]

Con ese valor, la derivada sigue siendo suficientemente cercana a la ideal para la banda de interes, pero se evita una amplificacion excesiva del ruido.

\subsubsection{Paso 6: valores finales para el bloque}

Reuniendo todo lo anterior:

\begin{itemize}[leftmargin=1.2cm]
\item del PID ideal se obtuvo $K_c=50$, $\tau_I=0.05$ y $\tau_D=0.2$;
\item del bloque se sabe que $P=K_c$, $I=\tau_I$ y $D=\tau_D$;
\item para el filtro derivativo se escoge $N=10$.
\end{itemize}

Por tanto, los valores finales para el bloque de Simulink son:
\begin{equation}
\boxed{
P=50,\qquad I=0.05,\qquad D=0.2,\qquad N=10
}
\label{eq:valores_bloque}
\end{equation}

\subsubsection{Paso 7: verificacion sustituyendo en el bloque}

Si se reemplazan estos valores en \eqref{eq:bloque_pid}, el bloque queda:
\begin{equation}
C_b(s)=50\left(1+\frac{1}{0.05\,s}+0.2\frac{10}{1+10/s}\right).
\end{equation}

Usando \eqref{eq:term_deriv_block_2}:
\begin{equation}
C_b(s)=50\left(1+\frac{1}{0.05\,s}+0.2\frac{10s}{s+10}\right).
\end{equation}

Como
\[
0.2\times 10 = 2,
\]
tambien puede escribirse como:
\begin{equation}
C_b(s)=50\left(1+\frac{1}{0.05\,s}+2\frac{s}{s+10}\right).
\label{eq:block_final_expanded}
\end{equation}

Si $N=10$ es suficientemente grande para la banda de trabajo, entonces:
\[
2\frac{s}{s+10}\approx 0.2\,s,
\]
y el bloque se comporta de forma muy cercana al controlador ideal
\[
50\left(1+\frac{1}{0.05\,s}+0.2\,s\right).
\]

\section{5. Evaluacion del desempe\~no}

El enunciado pide evaluar el controlador tanto sobre el modelo matematico aproximado como sobre el modelo realista. Deben estudiarse al menos:

\begin{itemize}[leftmargin=1.2cm]
\item rechazo a perturbaciones;
\item seguimiento de una referencia escalon variable;
\item tiempo de subida;
\item tiempo de establecimiento;
\item sobreimpulso;
\item error en estado estacionario.
\end{itemize}

\subsection{Diferencia entre regulacion y seguimiento}

Aunque el mismo controlador se usa en ambos escenarios, el significado fisico de las pruebas no es el mismo.

\subsubsection{Regulacion con setpoint fijo}

En la aplicacion real del laboratorio, la condicion natural es:
\[
r(t)=0.
\]

Eso significa que la plataforma del sensor no esta tratando de seguir una trayectoria variable, sino de permanecer inmovil mientras la base introduce perturbaciones. Este es un problema de regulacion con rechazo de perturbaciones.

En este escenario, las preguntas relevantes son:

\begin{itemize}[leftmargin=1.2cm]
\item cual es la desviacion maxima de la plataforma cuando aparece la perturbacion;
\item que tan rapido retorna cerca del equilibrio;
\item cuanta vibracion residual queda despues del transitorio;
\item si se presenta o no un error estacionario apreciable.
\end{itemize}

Este es el caso mas importante para el proyecto, porque el sensor ultrasónico necesita una base lo mas quieta posible para no contaminar la medicion.

\subsubsection{Seguimiento de referencia escalon}

El enunciado tambien pide analizar seguimiento de una referencia escalon variable. En ese caso la prueba ya no es una perturbacion sobre una referencia fija, sino un cambio de referencia:
\[
r(t)=r_0\,1(t).
\]

Aqui las metricas clasicas
\[
t_r,\qquad t_s,\qquad M_p,\qquad e_{ss}
\]
tienen una interpretacion directa. El objetivo es verificar si la plataforma puede alcanzar un nuevo valor de referencia con rapidez, poco sobreimpulso y error final pequeno.

La diferencia conceptual es entonces:

\begin{itemize}[leftmargin=1.2cm]
\item en regulacion, interesa mantener la salida cercana a cero a pesar de la perturbacion;
\item en seguimiento, interesa llevar la salida hacia un nuevo valor deseado.
\end{itemize}

\subsection{Lazo cerrado Simscape para rechazo de perturbaciones}

\figplaceholder{Respuesta en lazo cerrado del modelo realista en Simscape con el controlador IMC-PID para setpoint fijo y perturbacion de base.}{fig:lazo_cerrado_simscape}{Insertar aqui la grafica del modelo realista en lazo cerrado cuando la referencia es fija y la base vibra.}

En esta figura se debe analizar si el controlador reduce la transmision de vibracion hacia la plataforma y si la plataforma del sensor permanece suficientemente estable para la medicion.

\subsection{Lazo cerrado para seguimiento de referencia}

\figplaceholder{Respuesta en lazo cerrado para seguimiento de una referencia escalon.}{fig:lazo_cerrado_tracking}{Insertar aqui la grafica de seguimiento de referencia escalon para analizar rapidez, sobreimpulso y error final.}

Esta prueba es complementaria. No representa el caso principal del laboratorio, pero si permite evaluar la capacidad de servo del controlador y reportar las metricas clasicas pedidas en el curso.

\subsection{Tabla de metricas}

\begin{table}[H]
\centering
\small
\begin{tabular}{lcccc}
\toprule
\textbf{Caso} & \textbf{$t_r$} & \textbf{$t_s$} & \textbf{$M_p$} & \textbf{$e_{ss}$} \\
\midrule
Lazo abierto teorico & -- & -- & -- & -- \\
Lazo abierto Simscape & -- & -- & -- & -- \\
Lazo cerrado Simscape, regulacion & -- & -- & -- & -- \\
Lazo cerrado Simscape, seguimiento & -- & -- & -- & -- \\
\bottomrule
\end{tabular}
\caption{Metricas de desempe\~no a completar con los resultados de simulacion.}
\label{tab:metricas}
\end{table}

La tabla anterior debe completarse con las metricas obtenidas de las simulaciones y sirve para respaldar cuantitativamente el analisis cualitativo de las figuras.

Para el caso de regulacion con perturbaciones, ademas de las metricas anteriores, conviene reportar tambien la desviacion maxima de la plataforma, el tiempo de recuperacion y una medida de vibracion residual. Esas cantidades son especialmente importantes porque el objetivo practico del sistema es aislar el sensor, no solo responder bien a una referencia artificial.

\section{Espacio para bono}

El enunciado indica que se puede obtener una bonificacion adicional si el proyecto incluye una visualizacion tridimensional del sistema. Por ello se deja esta seccion opcional.

\figplaceholder{Visualizacion tridimensional opcional del sistema.}{fig:bono_3d}{Insertar aqui la captura o figura de la visualizacion 3D si se desarrolla la bonificacion.}

\section{Uso de IA y apoyo}

Para la elaboracion del presente documento se utilizo apoyo de inteligencia artificial como herramienta de sintesis, redaccion tecnica y reorganizacion del contenido del proyecto. En particular, la IA se uso para ayudar a convertir los desarrollos del proyecto en un documento continuo, ordenado y explicativo.

Los calculos del modelo, la comparacion de expresiones, la seleccion de parametros y la interpretacion fisica principal fueron desarrollados manualmente y luego trasladados como insumo para la redaccion del documento. Por tanto, la IA se utilizo como apoyo de presentacion y estructuracion, no como reemplazo del analisis de control ni del planteamiento fisico.

Adicionalmente, algunas reducciones algebraicas y la solucion simbolica de ecuaciones se apoyaron en \texttt{SymPy} dentro de un \textit{Jupyter Notebook}. Como respaldo se adjunta el archivo:
\[
\texttt{apoyo\_sympy\_reduccion\_imc\_pid.ipynb}
\]
en el que se dejan escritas la planta nominal, la reduccion de tercer a segundo orden y las expresiones simbolicas del controlador IMC-PID.

Tambien se utilizo el apoyo de \texttt{Codex} para obtener explicaciones paso a paso sobre la implementacion del sistema en Simscape y sobre la estructuracion de la implementacion tridimensional correspondiente al bono opcional.

\section{Conclusiones}

Se propuso una plataforma antivibratoria activa para estabilizar un sensor de ultrasonido usado en pruebas de laboratorio sobre papas. El modelo fisico completo se obtuvo a partir de principios mecanicos y electricos. A partir de ese modelo se construyo una planta nominal de tercer orden, se analizo su estabilidad y se dise\~no un controlador IMC-PID usando un modelo reducido de segundo orden para sintonia.

La reduccion del modelo se justifico suponiendo que la dinamica electrica del actuador es mas rapida que la dinamica mecanica dominante. Esa simplificacion permitio obtener una ley PID interpretable sin perder de vista que la validacion final debe hacerse sobre la planta completa y sobre el modelo realista en Simscape.

Con los parametros mecanicos dados, y manteniendo $P=50$, el controlador final queda:
\[
C(s)=50\left(1+\frac{1}{0.05\,s}+0.2\,s\right),
\]
y su forma equivalente para el bloque PID de Simulink es:
\[
P=50,\qquad I=0.05,\qquad D=0.2,\qquad N=10.
\]

Finalmente, se distinguieron dos escenarios de evaluacion: regulacion con setpoint fijo para rechazo de perturbaciones, que es el caso central de la aplicacion del sensor, y seguimiento de referencia escalon, que se incluye para cumplir el requisito de desempeno del curso. El documento queda estructurado de forma consistente con el enunciado y listo para completar con las figuras y metricas finales de simulacion.

\end{document}
